\documentclass[10pt,a4paper]{amsart}
\usepackage[margin=19mm]{geometry}
\usepackage{amsmath,amssymb,mathtools}
\usepackage[scr=rsfs,cal=euler]{mathalfa}
\usepackage{xfrac}
\usepackage[unicode,hidelinks,bookmarks=false]{hyperref}
\newcommand{\ie}{i.e.\ }
\newcommand{\cf}{cf.\ }
\newcommand{\resp}{respectively\ }
\newcommand{\Subsection}{\subsection}
\providecommand{\nobreakdash}{\nobreak\hbox{-}}
\setlength{\emergencystretch}{2em}
\multlinegap0pt
\usepackage{amssymb}
\usepackage{xcolor}
\usepackage{braket}
\usepackage{bm}
\usepackage[normalem]{ulem}
\usepackage{mathtools}
\newtheorem{theorem}{Theorem}
\newtheorem{lemma}{Lemma}
\title{A general proof of integer R\'{e}nyi QNEC}
\author{Tanay Kibe, Pratik Roy}
\date{Source: arXiv:2605.15272v2 (CC BY 4.0)}
\begin{document}
\maketitle

\noindent\emph{Source and licence.} A general proof of integer R\'{e}nyi QNEC. Version arXiv:2605.15272v2 (CC BY 4.0), distributed by arXiv under the Creative Commons Attribution 4.0 International licence, http://creativecommons.org/licenses/by/4.0/, as recorded in the arXiv licence metadata for this exact version and shown on the abstract page https://arxiv.org/abs/2605.15272v2. Full licence notice: \texttt{licenses/arxiv-2605.15272-CC-BY-4.0.txt}.\par\smallskip

\noindent\textit{Editorial scope.} Counted unit: the article's lemma lem:modular-Ln-group (source0.tex lines 3279--3308). The article's complete original proof of it is delivered with it (source0.tex lines 3310--3431, one \texttt{proof} environment of 122 lines). No mathematics is written, repaired or re-derived: the statement and the proof are the article's own lines, in the article's own notation, and no retained line is substituted or re-flowed.  Retained uncounted as the source-local context the proof needs: lines 1024--1039; lines 3065--3074.  Nothing was omitted from any retained span, so the package carries no omission record.  No retained line contains a citation, so the wrapper carries no bibliography and no bibliography entry.  No retained line contains a figure environment, an image-inclusion command or drawing code, so no image is delivered and the package declares no asset dependency.  Editorial formatting only, in addition: an AMS \texttt{amsart} preamble that restates the article's own declarations and macros used by the retained lines, namely the environments \texttt{theorem}, \texttt{lemma} on separate counters, together with the standard packages those lines need.  The retained environments print in this document as Theorem 1, Lemma 1 in order of appearance, whereas the article prints the same items with the numbering of its own full document.  The complete native source of the article as distributed in the arXiv e-print for this exact version is bundled unchanged as \texttt{sources/200642/source0.tex}.
\begin{theorem}[Extension of maps to $L^p$ spaces]\label{thm:extension-Lp}
Let \(M\) be a \(\sigma\)-finite von Neumann algebra, let
\(\omega\in M_*^+\) be faithful and normal, and let
\(T:M\to M\) be a normal positive unital map satisfying
\[
        \omega\circ T=\omega.
\]
For \(1\leq p<\infty\), define the map on the dense subspace
\[
    T^{(p)}: j_p(M)\longmapsto j_p(T(M)).
\]
Then $T^{(p)}$ extends uniquely to a positive contraction
\[
        T^{(p)}:L^p(M)\longrightarrow L^p(M)\,.
\]
\end{theorem}

\begin{lemma}[Sandwich continuity]\label{lem:sandwich-continuity}
Let \(h\in L^1(M)_+\), let \(1\le p<\infty\), and let
\((a_i)_i\subset M\) be uniformly bounded with \(a_i\to a\) strongly.  Then
\[
        h^{1/(2p)}a_i h^{1/(2p)}
        \longrightarrow
        h^{1/(2p)}a h^{1/(2p)}
        \qquad\text{in }L^p(M).
\]
\end{lemma}

\begin{lemma}[Strongly continuous modular isometry group on \(L^n(M)\)]
\label{lem:modular-Ln-group}
Define $\sigma_s:= \sigma_s^\omega$.
For every \(s\in\mathbb R\), the map
\[
        j_n(a)\longmapsto j_n(\sigma_s(a)),
        \qquad a\in M,
\]
extends uniquely to a positive surjective isometry
\[
        \Sigma_s^{(n)}:L^n(M)\longrightarrow L^n(M).
\]
Moreover,
\begin{equation}
\label{eq:Sigma-group-property}
        \Sigma_{s+t}^{(n)}
        =
        \Sigma_s^{(n)}\Sigma_t^{(n)},
        \qquad
        \Sigma_0^{(n)}=I,
        \qquad s,t\in\mathbb R,
\end{equation}
and the map
\[
        \mathbb R\ni s\longmapsto \Sigma_s^{(n)}x\in L^n(M)
\]
is norm-continuous for every \(x\in L^n(M)\).  Thus
\((\Sigma_s^{(n)})_{s\in\mathbb R}\) is a strongly continuous group of
positive isometries on \(L^n(M)\).
\end{lemma}

\begin{proof}
For each \(s\in\mathbb R\), the modular automorphism
\(\sigma_s:M\to M\) is normal, unital, positive, and
\(\omega\)-preserving.  Hence Theorem~\ref{thm:extension-Lp} shows that the
map
\[
        j_n(a)\longmapsto j_n(\sigma_s(a))
\]
extends uniquely to a positive contraction
\[
        \Sigma_s^{(n)}:L^n(M)\longrightarrow L^n(M).
\]

We first verify the group property.  For \(a\in M\) and \(s,t\in\mathbb R\),
\[
\begin{aligned}
        \Sigma_s^{(n)}\Sigma_t^{(n)}j_n(a)=
        j_n(\sigma_s(\sigma_t(a)))  =
        j_n(\sigma_{s+t}(a)) =
        \Sigma_{s+t}^{(n)}j_n(a).
\end{aligned}
\]
Since \(j_n(M)\) is norm-dense in \(L^n(M)\) and both sides are bounded,
the equality extends to all of \(L^n(M)\).  Similarly,
\(\Sigma_0^{(n)}=I\).  In particular,
\[
        \Sigma_{-s}^{(n)}\Sigma_s^{(n)}
        =
        \Sigma_s^{(n)}\Sigma_{-s}^{(n)}
        =
        I,
\]
so \(\Sigma_s^{(n)}\) is bijective with inverse
\(\Sigma_{-s}^{(n)}\).

Since both \(\Sigma_s^{(n)}\) and \(\Sigma_{-s}^{(n)}\) are contractions,
for every \(x\in L^n(M)\),
\[
        \|x\|_n
        =
        \|\Sigma_{-s}^{(n)}\Sigma_s^{(n)}x\|_n
        \leq
        \|\Sigma_s^{(n)}x\|_n
        \leq
        \|x\|_n.
\]
Consequently,
\[
        \|\Sigma_s^{(n)}x\|_n=\|x\|_n,
\]
and \(\Sigma_s^{(n)}\) is an isometry.

It remains to prove strong continuity.  We first prove continuity at
\(s=0\) on the dense subspace \(j_n(M)\).  Let \(a\in M\).  The modular
automorphism group is point strong-\(*\) continuous, and
\[
        \|\sigma_s(a)\|=\|a\|,
        \qquad s\in\mathbb R.
\]
Therefore,
\[
        \sigma_s(a)\longrightarrow a
        \qquad\text{strongly as }s\to0,
\]
with the family \((\sigma_s(a))_{s\in\mathbb R}\) uniformly bounded.
Lemma~\ref{lem:sandwich-continuity} then gives
\begin{equation}
\label{eq:Sigma-cont-core}
\begin{aligned}
        \|\Sigma_s^{(n)}j_n(a)-j_n(a)\|_n
        &=
        \|j_n(\sigma_s(a))-j_n(a)\|_n
        \longrightarrow0
        \qquad (s\to0).
\end{aligned}
\end{equation}

Now let \(x\in L^n(M)\) and let \(\varepsilon>0\).  Choose
\(y\in j_n(M)\) such that
\[
        \|x-y\|_n<\varepsilon.
\]
Since \(\Sigma_s^{(n)}\) is an isometry,
\[
\begin{aligned}
        \|\Sigma_s^{(n)}x-x\|_n
        &\leq
        \|\Sigma_s^{(n)}(x-y)\|_n
        +
        \|\Sigma_s^{(n)}y-y\|_n
        +
        \|y-x\|_n                                    \\
        &\leq
        2\varepsilon
        +
        \|\Sigma_s^{(n)}y-y\|_n.
\end{aligned}
\]
By \eqref{eq:Sigma-cont-core}, the last term tends to zero as \(s\to0\).
Since \(\varepsilon>0\) is arbitrary,
\[
        \Sigma_s^{(n)}x\longrightarrow x
        \qquad\text{in }L^n(M)
\]
as \(s\to0\).

Finally, for any \(s_0\in\mathbb R\), the group property and isometry give
\[
\begin{aligned}
        \|\Sigma_s^{(n)}x-\Sigma_{s_0}^{(n)}x\|_n
        &=
        \left\|
        \Sigma_{s_0}^{(n)}
        \bigl(\Sigma_{s-s_0}^{(n)}x-x\bigr)
        \right\|_n                                   \\
        &=
        \|\Sigma_{s-s_0}^{(n)}x-x\|_n
        \longrightarrow0
\end{aligned}
\]
as \(s\to s_0\).  This proves strong continuity on \(\mathbb R\).
\end{proof}

\end{document}
